\documentclass{article}

% 1) Seitenränder mit geometry festlegen
\usepackage[a4paper,
  left=25mm,      % Linker Rand: 25mm
  top=25mm,       % Oberer Rand: 25mm
  bottom=25mm,    % Unterer Rand: 25mm
  right=50mm      % Rechter Rand: 50mm (5cm)
]{geometry}

% 2) Pakete für das Zwei-Spalten-Layout und Farben
\usepackage{paracol}
\usepackage{xcolor}
\usepackage{tikz}

% 3) Fix für \tightlist (Pandoc-Standard)
\providecommand{\tightlist}{%
  \setlength{\itemsep}{0pt}\setlength{\parskip}{0pt}%
}

% 4) Schriftart festlegen (Open Sans)
\usepackage{fontspec}
\setmainfont{Open Sans}

\begin{document}

% 5) Start des Zwei-Spalten-Layouts
\begin{paracol}{2}

% 6) Linke Spalte: Haupttext
\switchcolumn[0]
% Quarto fügt hier den Haupttext ein
$body$

% 7) Rechte Spalte: Platz für zukünftigen Inhalt (z.B. Logo, Anschrift)
\switchcolumn
% Platzhalter für den rechten Rand
% Hier kannst du später Logo, Anschrift und Seitenzahlen hinzufügen

\end{paracol}

\end{document}